We study six parliamentary democracies chosen to span the space of electoral systems, seat magnitudes, and redistricting institutions (Table~\ref{tab:country-comparison}).

\begin{table}[h]
\centering
\caption{Redistricting systems in six parliamentary democracies.}
\label{tab:country-comparison}
\small
\begin{tabular}{llrrll}
\toprule
Country & System & Districts & Seats & Balance basis & Commission \\
\midrule
United Kingdom & FPTP & 543 & 543 & Registered electors & Independent \\
Canada & FPTP & 343 & 343 & Total population & Independent \\
Ireland & STV & 43 & 174 & Per seat & Independent \\
Malta & STV & 13 & 65 & Per seat & Parliament \\
New Zealand & MMP & 65 & 65 & Electoral pop. & Independent \\
Germany & MMP & 299 & 299+ & German citizens & Advisory \\
\bottomrule
\end{tabular}
\end{table}

\subsection{United Kingdom}

The UK uses single-member constituencies (``seats'') for the House of Commons elected by first-past-the-post plurality. Four independent Boundary Commissions (England, Scotland, Wales, Northern Ireland) draw constituency boundaries using registered electorate --- not total population --- as the balance criterion. The Parliamentary Constituencies Act 2020 sets the electoral quota at the national electorate divided by 650 seats, with 5\% tolerance. Constituencies must respect local authority boundaries where possible. Our analysis covers England (543 constituencies) where the largest commission operates and data are most available.

\subsection{Canada}

Canada's 343 federal ridings are drawn by provincial Electoral Boundary Commissions appointed by judges. The process uses total population with 25\% deviation from provincial quotas permitted; communities of interest (including Indigenous communities) are a named criterion. Unlike the UK, no single national commission exists; each of Canada's 10 provinces has its own commission. This decentralized structure creates variation in commission independence and methodology.

\subsection{Ireland}

Ireland's Dáil Éireann uses single transferable vote (STV) in 43 multi-member constituencies returning 3--5 seats each (174 total seats as of 2024). The Electoral Commission draws boundaries using total population with a per-seat target (national population $\div$ 174). Because constituencies have different seat counts, balance is measured \emph{per seat}: a 5-seat constituency must have approximately $5 \times \text{ideal\_per\_seat}$ population, and a 3-seat constituency must have $3 \times \text{ideal\_per\_seat}$. County boundaries must be preserved where possible.

\subsection{Malta}

Malta's House of Representatives uses STV in 13 districts, each returning 5 seats (65 base seats; additional seats may be added for proportionality). All constituencies are constitutionally fixed to return exactly 5 seats, making Malta a clean test case for uniform-magnitude multi-member redistricting. The parliament draws boundaries with limited independent oversight.

\subsection{New Zealand}

New Zealand uses MMP, with 65 general electorates (single-member, FPTP) and 7 Māori electorates drawn separately, plus list seats allocated for proportionality (120 total). The Representation Commission draws both general and Māori electorate boundaries using the \emph{general electoral population} (a Census-derived count excluding those enrolled on the Māori roll). This paper covers general electorates only.

\subsection{Germany}

Germany's 299 Bundeswahlkreise (federal electoral districts) are single-member constituencies using FPTP; list seats are allocated to ensure MMP proportionality. The Wahlkreiskommission is advisory, with the Bundestag retaining final authority over boundaries. The population basis is \emph{German citizens} only --- approximately 10 million non-citizen residents are excluded. Boundaries must not divide Landkreise (rural counties) or kreisfreie Städte (independent cities) where possible, with 25\% tolerance from the national average.

\subsection{Data}

% TODO: Describe census-equivalent data sources for each country.
% UK: Office for National Statistics census output areas (2021)
% Canada: Statistics Canada dissemination areas (2021)
% Ireland: Central Statistics Office small areas (2022)
% Malta: National Statistics Office enumeration districts (2021)
% NZ: Stats NZ meshblocks (2023)
% Germany: Destatis Gemeindeverbände (2021)
% All data is publicly available; geographic unit GEOIDs encoded as
% {ISO-2-country}{4-char-unit-code} in RPLAN format.
